\booksection{บทที่ 20: โหมดนักพัฒนาและสิทธิ์ระดับสูง}

\chapterintro{บทนี้เขียนไว้สำหรับคนที่อยากเจาะลึก “ฟังก์ชันที่ซ่อนอยู่” ของมือถือ โหมดนักพัฒนา สิทธิ์ Root กุญแจ BL ... แนวคิดพวกนี้ฟังดูเหมือนของสายเกจ แต่แท้จริงไม่ได้ห่างจากผู้ใช้ทั่วไปเลย}

\booksubsection{หนึ่ง โหมดนักพัฒนา: ปลดล็อกการตั้งค่าที่ซ่อนอยู่}

โหมดนักพัฒนาคือเมนูลับที่มากับระบบมือถือ ข้างในเต็มไปด้วยฟังก์ชันที่ผู้ใช้ทั่วไปไม่ค่อยได้ใช้ แต่นักพัฒนาจำเป็นต้องใช้

\subsubsection*{วิธีเปิด:}

\begin{itemize}
  \item \textbf{Android}: เข้า การตั้งค่า > เกี่ยวกับโทรศัพท์ > กด "หมายเลขรุ่น" ต่อเนื่อง 7 ครั้ง จนขึ้นข้อความว่า "คุณอยู่ในโหมดนักพัฒนาแล้ว"
  \item \textbf{ฮาร์โมนีโอเอส}: การตั้งค่า > เกี่ยวกับโทรศัพท์ > กด "เวอร์ชัน HarmonyOS" ต่อเนื่อง 7 ครั้ง
  \item \textbf{iOS}: ผู้ใช้ทั่วไปเปิดโหมดนักพัฒนาแบบเต็มรูปแบบไม่ได้ ต้องมีบัญชี Apple Developer แบบเสียเงิน
\end{itemize}

\subsubsection*{รวมฟังก์ชันที่ใช้บ่อย:}

\begin{itemize}
  \item \textbf{การดีบัก USB (USB Debugging)}:สวิตช์ที่ต้องเปิดสำหรับเชื่อมต่อคอมพิวเตอร์เพื่อทำงานขั้นสูง (เช่น การแฟลชรอม, กู้คืนข้อมูล, คำสั่ง ADB) ดูรายละเอียดเพิ่มเติมได้ในส่วน "คำสั่ง ADB ที่ใช้งานได้จริง" ของบทนี้
  \item \textbf{แสดงการสัมผัสหน้าจอ (Show Taps)}:แสดงจุดสัมผัสบนหน้าจอ (จุดกลมเล็กๆ ตามนิ้ว) เหมาะสำหรับบันทึกหรือสาธิตวิธีการใช้งาน
\end{itemize}

\subsubsection*{ที่เกี่ยวข้องกับประสิทธิภาพและความลื่นไหล:}

\begin{itemize}
  \item \textbf{ปรับความเร็วแอนิเมชัน (สเกลแอนิเมชันหน้าต่าง / สเกลแอนิเมชันการเปลี่ยนผ่าน / สเกลระยะเวลาแอนิเมชันโปรแกรม)}:ตัวเลือกทั้งสามนี้ควบคุมระยะเวลาแอนิเมชันตอนเปิด App หรือสลับหน้าจอ ค่าเริ่มต้นคือ 1x ปรับเป็น 0.5x จะทำให้โทรศัพท์ "ดูเหมือน" เร็วขึ้น (ความเร็วการทำงานจริงของโปรแกรมไม่เปลี่ยน เป็นแค่แอนิเมชันที่เร็วขึ้น) หากตั้งเป็นปิดจะข้ามไปยังหน้าจอใหม่ทันทีโดยไม่มีแอนิเมชัน ส่วนใหญ่ผู้ที่ชื่นชอบเทคโนโลยีมักจะตั้งไว้ที่ 0.5x เพื่อให้ยังคงความรู้สึกตอนเปลี่ยนหน้าจอแต่ดูเร็วขึ้น การปิดทั้งหมดจะทำให้การใช้งานดูไม่ลื่นไหล และอาจทำให้มองข้ามอาการกระตุกจริงของโปรแกรม จนทำให้คุณเข้าใจผิดว่า "ลื่นไหลมาก"
  \item \textbf{บังคับ GPU เรนเดอร์ (Force GPU Rendering)}:บังคับให้กราฟิก 2D ทั้งหมดประมวลผลผ่าน GPU แทน CPU ช่วยเพิ่มความลื่นไหลตอนเลื่อนหน้าจอในเครื่องรุ่นเก่าหรือบาง App ที่ปรับแต่งไม่ดี แต่ก็เพิ่มภาระให้ GPU และกินไฟมากขึ้น โทรศัพท์รุ่นใหม่ใช้ GPU เรนเดอร์เป็นค่าเริ่มต้นอยู่แล้ว สวิตช์นี้จึงมีผลจำกัดกับเครื่องรุ่นกลางถึงระดับเรือธง ถือเป็นตัวเลือกที่เหลือมาจากอดีต ผู้ใช้ทั่วไปแค่ตั้งค่าเริ่มต้นไว้ก็พอ
  \item \textbf{แสดงการวิเคราะห์โหมดการเรนเดอร์ GPU (Profile GPU Rendering)}:แสดงเวลาที่ใช้ในการเรนเดอร์แต่ละเฟรมแบบเรียลไทม์บนหน้าจอเป็นกราฟแท่ง เส้นสีเขียวแนวนอนแทน 16.6ms (ค่าเกณฑ์เวลาเฟรมที่ 60fps) หากแท่งกราฟเกินเส้นสีเขียวแสดงว่ามีอาการดรอปเฟรม นี่เป็นเครื่องมือสำหรับนักพัฒนาเพื่อไล่หาสาเหตุ App กระตุก ผู้ใช้ทั่วไปก็สามารถใช้เปรียบเทียบได้ว่า App ไหนปรับแต่งได้ไม่ดี — เช่น ถ้า App ใด App หนึ่งตอนเลื่อนหน้าจอกราฟมักจะพุ่งสูงกว่าเส้นแดงเสมอ ก็แทบจะสรุปได้ว่าเป็นความผิดของ App นั้น
\end{itemize}

\subsubsection*{ที่เกี่ยวข้องกับการทำงานเบื้องหลังและแบตเตอรี่:}

\begin{itemize}
  \item \textbf{จำกัดโปรเซสเบื้องหลัง}:สามารถตั้งค่าเป็น "จำกัดมาตรฐาน" "สูงสุด 1 โปรเซส" "สูงสุด 2 โปรเซส" เป็นต้น ดูเผินๆ การจำกัดโปรเซสเบื้องหลังจะช่วยประหยัดแบตเตอรี่ แต่จริงๆ แล้วระบบจัดการเบื้องหลังของ Android ทำมาได้ดีแล้ว การบังคับจำกัดอาจทำให้ App ต้องโหลดใหม่บ่อยๆ (การเริ่มต้นเย็นกลับกลายเป็นการกินไฟมากกว่าและประสบการณ์การใช้งานก็แย่ลง) ตัวเลือกนี้เหมาะกับการดีบักมากกว่า ไม่แนะนำให้ผู้ใช้ทั่วไปเปิดใช้งานระยะยาว ตั้งค่าเป็น "จำกัดมาตรฐาน" ไว้ก็พอ
  \item \textbf{ไม่เก็บกิจกรรม}:ทันทีที่ออกจากหน้าจอใดๆ จะทำลายมันทิ้งเลย ไม่แนะนำอย่างยิ่ง เพราะจะทำให้ App ค้างจอขาวและโหลดใหม่บ่อยๆ ใช้สำหรับทดสอบตรรกะการกู้คืนของ App เท่านั้น
\end{itemize}

\begin{tcbwarn}{ข้อควรระวัง}
การตั้งค่าในโหมดนักพัฒนาอาจกระทบความเสถียรของระบบ อย่าแก้ตัวเลือกที่ไม่เข้าใจ หลังเปิด USB Debugging ตอนต่อคอมพิวเตอร์ต้องระวัง อย่าเชื่อมต่อกับคอมที่ไม่น่าเชื่อถือ

\end{tcbwarn}

\booksubsection{ปลดล็อก OEM: “สวิตช์สำหรับแฟลชรอม” ในตัวเลือกนักพัฒนา}

ในตัวเลือกนักพัฒนามีสวิตช์ที่ไม่ค่อยสะดุดตาแต่สำคัญมาก — “ปลดล็อก OEM” (บางแบรนด์เรียก “อนุญาตการปลดล็อก OEM” หรือ “ปลดล็อกอุปกรณ์”) สวิตช์นี้เองไม่ได้ปลดล็อก BL ให้โดยตรง แต่มันคือ “ใบอนุญาตเบื้องต้น”: ต้องเปิดมันก่อน ผู้ผลิตถึงจะอนุญาตให้คุณปลดล็อก BL ผ่านคำสั่ง fastboot หรือเครื่องมือของแบรนด์ได้

\begin{itemize}
  \item \textbf{หน้าที่ของสวิตช์}:เทียบเท่ากับการประกาศต่อผู้ผลิตว่า "ฉันยอมสละสิทธิ์การรับประกัน และยอมรับความเสี่ยงเอง โปรดอนุญาตให้ฉันปลดล็อก BL" หากสวิตช์นี้ปิดอยู่ คำสั่งปลดล็อก BL ใดๆ ก็จะถูกปฏิเสธทันที
  \item \textbf{ความสัมพันธ์กับ BL Lock}:ปลดล็อก OEM คือ "การขอสิทธิ์" ส่วนการปลดล็อก BL คือ "การดำเนินการจริง" อันแรกทำได้แค่แตะในการตั้งค่าของโทรศัพท์ ส่วนอันหลังต้องไปขอรหัสปลดล็อกจากเว็บไซต์ผู้ผลิต แล้วใช้คอมพิวเตอร์รันคำสั่ง fastboot (เช่น\texttt{fastboot flashing unlock})ทั้งสองอย่างนี้ขาดไม่ได้
  \item \textbf{การรีเซ็ตเป็นค่าจากโรงงานจะรีเซ็ตหรือไม่}:จะรีเซ็ต สวิตช์ "ปลดล็อก OEM" จะกลับไปอยู่ในสถานะปิดหลังจากรีเซ็ตเป็นค่าจากโรงงาน แต่สถานะการปลดล็อก BL นั้นเป็นการเปลี่ยนแปลงในระดับฮาร์ดแวร์ เมื่อปลดล็อกแล้วจะไม่ล็อกกลับแม้จะรีเซ็ตเป็นค่าจากโรงงาน — ต้องสั่ง\texttt{fastboot flashing lock}ด้วยตนเองเท่านั้นถึงจะล็อกกลับได้
  \item \textbf{เมื่อไรควรเปิด}:เว้นแต่คุณแน่ใจว่าจะแฟลชรอมหรือ Root ไม่เช่นนั้นให้ปิดไว้ มันเป็นสวิตช์ที่แยกจากกันกับ "USB Debugging" ตอนแฟลชรอมมักจะต้องเปิดทั้งสองตัวพร้อมกัน
\end{itemize}

\begin{tcbinfo}{เกร็ดความรู้: ทำไมผู้ผลิตถึงตั้งสวิตช์นี้ไว้}
ในยุคแรกของ Android ไม่มีสวิตช์นี้ มัลแวร์สามารถแอบแฟลชระบบได้เพียงแค่ได้สิทธิ์เข้าถึง Google จึงกำหนดให้ผู้ผลิต Android ทุกรายต้องใส่สวิตช์ “ปลดล็อก OEM” ในตัวเลือกนักพัฒนา เพื่อเป็นการยืนยัน “ความยินยอมโดยรู้ตัวของผู้ใช้” นี่คือเหตุผลที่แม้คุณจะปลดล็อก BL แล้ว แต่ถ้าสวิตช์ปลดล็อก OEM ถูกปิด การแฟลชรอมใหม่ก็จะล้มเหลว

\end{tcbinfo}

\booksubsection{BL Lock (Bootloader Lock): ประตูบานแรกของการแฟลชรอม}

Bootloader คือโปรแกรมที่โหลดระบบปฏิบัติการขึ้นมาตอนมือถือเปิดเครื่อง BL Lock คือกลอนที่ผู้ผลิตใส่ไว้เพื่อความปลอดภัย กันผู้ใช้แฟลชระบบที่ไม่ใช่ของทางการ

\subsubsection*{หน้าที่ของ BL Lock:}

\begin{itemize}
  \item \textbf{การปกป้องความปลอดภัย}: ป้องกันมัลแวร์แอบแฟลชเข้าระบบ
  \item \textbf{การกันขโมย}: แม้มือถือจะถูกแฟลชรอม ถ้าไม่มีรหัสผ่านบัญชีเดิมก็ใช้ไม่ได้
  \item \textbf{ข้อจำกัดการรับประกัน}: ปลดล็อก BL แล้วมักแปลว่าเสียการรับประกันจากทางการ
\end{itemize}

\subsubsection*{นโยบาย BL Lock ของแต่ละแบรนด์:}

\begin{itemize}
  \item \textbf{Xiaomi}: ปลดล็อกได้ ต้องผูกบัญชีรอช่วงเวลา (7-30 วัน)
  \item \textbf{Huawei}: หลังปี 2018 หยุดแจกจ่ายรหัสปลดล็อก แทบไม่มีทางปลดล็อกอย่างเป็นทางการ
  \item \textbf{OnePlus}: ปลดล็อกได้ ค่อนข้างหลวม
  \item \textbf{OPPO/vivo}: บางรุ่นปลดล็อกได้ ต้องยื่นสมัคร
  \item \textbf{Samsung}: ปลดล็อกได้ แต่จะทริกเกอร์ Knox ตัดสลัก เสียฟังก์ชันระดับสูงอย่าง Samsung Pay
  \item \textbf{Apple}: อุปกรณ์ iOS ไม่มีแนวคิด BL Lock แบบดั้งเดิม แต่มีกลไกความปลอดภัยที่คล้ายกัน
\end{itemize}

\booksubsection{คำสั่ง ADB ที่ใช้งานได้จริง: ควบคุมโทรศัพท์ผ่านคอมพิวเตอร์}

ADB ย่อมาจาก Android Debug Bridge เป็นเครื่องมือคอมมานด์ไลน์ที่ Google จัดทำอย่างเป็นทางการ ช่วยให้คอมพิวเตอร์สื่อสารกับโทรศัพท์ Android ผ่าน USB หรือ Wi-Fi ได้ ตัวมันเองไม่ต้องใช้ Root ขอแค่เปิด “USB Debugging” ในโทรศัพท์ก็ใช้งานได้แล้ว สามารถเข้าใจได้ว่ามันคือ “รีโมตควบคุมที่ใช้คอมพิวเตอร์สั่งงานโทรศัพท์” — ไม่ว่าจะแฟลชรอม ติดตั้ง App ส่งไฟล์ หรือเช็ค Log ก็ต้องพึ่งมันทั้งนั้น

\subsubsection*{วิธีเชื่อมต่อผ่าน USB:}

\begin{enumerate}
  \item ฝั่งโทรศัพท์: เข้าไปที่ตัวเลือกนักพัฒนา เปิด "USB Debugging"
  \item ฝั่งคอมพิวเตอร์: ดาวน์โหลด Platform Tools จากเว็บไซต์ทางการของ Google (รวม\texttt{adb}และ\texttt{fastboot}ไฟล์ปฏิบัติการสองไฟล์) แตกไฟล์ไปยังโฟลเดอร์ใดก็ได้
  \item เชื่อมต่อมือถือกับคอมพิวเตอร์ด้วยสายเคเบิล บนมือถือจะมีหน้าต่างป๊อปอัป "อนุญาต USB debugging หรือไม่" ให้ติ๊ก "อนุญาตเสมอ" แล้วกดตกลง
  \item ใน Command Line ของคอมพิวเตอร์ (PowerShell หรือ cmd ของ Windows) ให้รัน\texttt{adb devices}ถ้าผลลัพธ์ขึ้นประมาณ\texttt{List of devices attached /\allowbreak  xxxxxxxx device}ก็แปลว่าเชื่อมต่อสำเร็จแล้ว
\end{enumerate}

\begin{tcbwarn}{อย่ามองข้ามไดรเวอร์และหน้าต่างป๊อปอัปการอนุญาต}
ผู้ใช้ Windows ถ้า\texttt{adb devices}ขึ้นว่า “unauthorized” หรือไม่แสดงรายการอุปกรณ์เลย ส่วนใหญ่เป็นเพราะติดตั้งไดรเวอร์ไม่ดีหรือลืมกด “อนุญาต” บนมือถือ แบรนด์อย่าง Samsung/ Xiaomi สามารถติดตั้ง “USB ไดรเวอร์” หรือ “โปรแกรมช่วยเหลือมือถือ” ของแต่ละแบรนด์เพื่อเติมไดรเวอร์ให้ครบได้ ส่วน Mac และ Linux มักจะไม่ต้องติดตั้งไดรเวอร์ ถ้าหน้าต่างป๊อปอัปบนมือถือถูกปิดไปโดยไม่ได้ตั้งใจ แค่ถอดเสียบสายเคเบิลใหม่ก็จะขึ้นมาอีกครั้ง

\end{tcbwarn}

\subsubsection*{คำสั่ง ADB ที่ใช้บ่อย:}

\begin{itemize}
  \item \texttt{adb install app.\allowbreak apk}:ติดตั้ง APK จากคอมพิวเตอร์ลงมือถือ (\texttt{-\allowbreak r}หมายถึงติดตั้งทับและเก็บข้อมูลเดิมไว้; ติดตั้ง\texttt{-\allowbreak s}ลง SD card)
  \item \texttt{adb uninstall com.\allowbreak xxx.\allowbreak xxx}:ถอนการติดตั้ง App ตามชื่อแพ็กเกจ
  \item \texttt{adb push พาธคอมพิวเตอร์ /\allowbreak sdcard/\allowbreak พาธปลายทาง}:ส่งไฟล์จากคอมพิวเตอร์ไปยังมือถือ
  \item \texttt{adb pull /\allowbreak sdcard/\allowbreak ไฟล์ พาธคอมพิวเตอร์}:ดึงไฟล์จากมือถือมายังคอมพิวเตอร์ (เร็วกว่าส่งให้ตัวเองบน WeChat แล้วค่อยดาวน์โหลดเยอะ)
  \item \texttt{adb shell}:เข้าสู่ Linux Shell ของมือถือ สามารถรันคำสั่ง\texttt{ls},\texttt{pm},\texttt{am}และอื่นๆ ได้โดยตรง
  \item \texttt{adb reboot bootloader}:รีสตาร์ทมือถือเข้าสู่โหมด Fastboot (ทางเข้าสำหรับการแฟลชรอม)
  \item \texttt{adb reboot recovery}:รีสตาร์ทเข้าสู่โหมด Recovery
  \item \texttt{adb shell pm list packages}:แสดงรายการชื่อแพ็กเกจของแอปที่ติดตั้งไว้ทั้งหมดบนมือถือ (สิ่งจำเป็นสำหรับการหาชื่อแพ็กเกจ) ใส่\texttt{-\allowbreak s}เพื่อดูเฉพาะแอประบบ ใส่\texttt{-\allowbreak 3}เพื่อดูเฉพาะแอปของบุคคลที่สาม
  \item \texttt{adb shell pm uninstall -\allowbreak -\allowbreak user 0 ชื่อแพ็กเกจ}:ถอนการติดตั้งแอปที่ติดมากับเครื่องสำหรับผู้ใช้ปัจจุบัน โดยไม่ต้อง Root ดูรายละเอียดเพิ่มเติมได้ในส่วนปฏิบัติจริงด้านล่าง
\end{itemize}

\subsubsection*{การถอนการติดตั้งแอปที่ติดมากับเครื่องโดยไม่ต้อง Root ตามจริง:}

นี่เป็นหนึ่งในฟังก์ชันที่ใช้งานได้จริงที่สุดของ ADB: การใช้คำสั่ง\texttt{pm uninstall -\allowbreak -\allowbreak user 0}สามารถถอนการติดตั้งแอปที่ติดมากับเครื่องสำหรับ “ผู้ใช้ปัจจุบัน” ได้โดยไม่ต้อง Root และไม่กระทบพาร์ติชันระบบ แอปที่ถอนการติดตั้งไปจะ “หายไป” จากระบบ แต่จริงๆ แล้วยังอยู่ในพาร์ติชันระบบ เมื่อรีเซ็ตเป็นค่าโรงงานจะกลับมาอีกครั้ง

\begin{enumerate}
  \item ก่อนอื่นใช้\texttt{adb shell pm list packages -\allowbreak s}เพื่อแสดงรายการแอประบบทั้งหมด แล้วหาชื่อแพ็กเกจของแอปที่คุณต้องการถอนการติดตั้ง ตัวอย่างเช่น เบราว์เซอร์ของ Xiaomi คือ\texttt{com.\allowbreak android.\allowbreak browser}ส่วน App Store ของ Huawei คือ\texttt{com.\allowbreak vmall.\allowbreak client}
  \item ไม่แน่ใจว่าชื่อแพ็กเกจไหนคือแอปไหน? สามารถใช้คอมพิวเตอร์รัน\texttt{adb shell pm list packages | findstr คำค้นหา}เพื่อค้นหา หรือติดตั้งแอปของบุคคลที่สามอย่าง "Package Viewer" เพื่อดูชื่อแพ็กเกจของแต่ละแอปได้โดยตรงบนมือถือ
  \item รันการถอนการติดตั้ง:\texttt{adb shell pm uninstall -\allowbreak -\allowbreak user 0 ชื่อแพ็กเกจ}เช่น\texttt{adb shell pm uninstall -\allowbreak -\allowbreak user 0 com.\allowbreak android.\allowbreak browser}ถ้าส่งค่ากลับมาเป็น\texttt{Success}ก็แปลว่าสำเร็จ
  \item อยากกู้คืนต้องทำยังไง? มี 2 วิธี: (1) รีเซ็ตเป็นค่าโรงงาน (ทำลายล้างที่สุดแต่ยุ่งยาก); (2) ติดตั้งใหม่:\texttt{adb shell cmd package install-\allowbreak existing ชื่อแพ็กเกจ}
\end{enumerate}

\textbf{อย่าลบแพ็กเกจหลักของระบบแบบสุ่มสี่สุ่มห้า}\par

แม้ว่า “pm uninstall --user 0” จะเป็นวิธีที่ปลอดภัยโดยไม่ต้อง Root แต่การลบแอประบบที่สำคัญแบบสุ่มสี่สุ่มห้า (เช่น\texttt{com.\allowbreak android.\allowbreak systemui}แถบสถานะ,\texttt{com.\allowbreak android.\allowbreak phone}โทรศัพท์,\texttt{com.\allowbreak android.\allowbreak settings}การตั้งค่า) อาจทำให้ระบบพัง หน้าจอดำ เปิดเครื่องไม่ติด และปัญหาร้ายแรงอื่นๆ ก่อนลบควร Google ดูก่อนว่า “ชื่อแพ็กเกจนี้ลบได้ไหม” ลบน้อยกว่าที่ควรลบยังดีกว่าลบผิด แนะนำให้ลบทีละหนึ่งหรือสองแอป แล้วยืนยันว่าไม่มีปัญหาก่อนค่อยดำเนินการต่อ

\begin{tcbinfo}{เกร็ดความรู้เล็กๆ: --user 0 หมายถึงอะไร}
หลังจาก Android 4.2 รองรับการใช้งานหลายผู้ใช้ โดยผู้ใช้แต่ละคนจะมี ID แยกจากกัน (เจ้าของเครื่องคือ 0 ผู้ใช้เพิ่มเติมจะเริ่มจาก 10 เป็นต้นไป)\texttt{-\allowbreak -\allowbreak user 0}หมายถึง “ถอนการติดตั้งสำหรับเจ้าของเครื่องปัจจุบันเท่านั้น” ไฟล์ APK ในพาร์ติชันระบบยังคงอยู่เหมือนเดิม เพียงแต่ไม่แสดงให้เจ้าของเครื่องเห็น นี่คือเหตุผลที่มันไม่ต้องใช้ Root เพราะไม่ได้แตะไฟล์ระบบ เพียงแค่เปลี่ยน “การกำหนดค่าผู้ใช้” นี่คือความแตกต่างที่ใหญ่ที่สุดระหว่างวิธีนี้กับการถอนการติดตั้งหลังจาก Root: การถอนการติดตั้งด้วย Root คือการลบออกจากระบบอย่างแท้จริง ผู้ใช้ทุกคนจะไม่มีแอปนั้นอีกต่อไป ส่วนการถอนการติดตั้งด้วย ADB เป็นเพียงการ “ซ่อนจากผู้ใช้ปัจจุบัน”

\end{tcbinfo}

\booksubsection{สิทธิ์ Root: “ผู้ดูแลระบบสูงสุด” ของ Android}

สิทธิ์ Root เทียบเท่ากับ Super User ใน Linux (มีสิทธิ์ควบคุมเต็มรูปแบบเหนือไดเรกทอรี / สามารถอ่านเขียนไฟล์ใดก็ได้) ถือเป็นสิทธิ์สูงสุดของระบบ Android ซึ่งมีความแตกต่างอย่างชัดเจนจาก Administrator ของ Windows: Windows ด้วยเหตุผลด้านความปลอดภัยได้จำกัดสิทธิ์ผู้ดูแลระบบไว้หลายประการ (สิทธิ์สูงสุดที่แท้จริงคือ TrustedInstaller) ดังนั้นแม้แต่ Administrator ก็ไม่สามารถเรียกใช้งานโดยตรงเหมือน Linux root ได้\texttt{sudo rm -\allowbreak rf /\allowbreak *}หลังจากได้รับสิทธิ์ Root แล้ว คุณสามารถแก้ไขไฟล์ระบบ ถอนการติดตั้งแอปที่มาพร้อมเครื่อง และใช้งานแอปขั้นสูงที่ต้องการ Root ได้

\subsubsection*{ข้อดีของ Root:}

\begin{itemize}
  \item : ถอนการติดตั้งแอปที่มาพร้อมเครื่องซึ่งปกติถอนไม่ได้
  \item : ใช้แอปฟังก์ชันขั้นสูง (เช่น Titanium Backup, Greenify)
  \item : แก้ไขอินเทอร์เฟซและฟังก์ชันของระบบ
  \item : ติดตั้ง Recovery บุคคลที่สามและ ROM กำหนดเอง
\end{itemize}

\subsubsection*{ความเสี่ยงของ Root:}

\begin{itemize}
  \item \textbf{เสียการรับประกัน}: ผู้ผลิตส่วนใหญ่ไม่รับประกันมือถือที่ผ่านการ Root
  \item \textbf{ความเสี่ยงด้านความปลอดภัย}: แอปที่ประสงค์ร้ายอาจได้สิทธิ์ Root ก่อผลเสียรุนแรง
  \item \textbf{ระบบไม่เสถียร}: ผิดพลาดเล็กน้อยอาจทำให้ระบบพัง
  \item \textbf{แอปธนาคารใช้ไม่ได้}: แอปธนาคารและแอปชำระเงินจำนวนมากตรวจพบ Root จะปฏิเสธการทำงานทันที
  \item \textbf{อัปเดต OTA ล้มเหลว}: อาจรับอัปเดตระบบอย่างเป็นทางการได้ไม่ปกติ
\end{itemize}

\textbf{คำเตือนร้ายแรง}\par

การ Root มีความเสี่ยง อาจทำให้มือถือกลายเป็นอิฐ (เปิดเครื่องไม่ได้) ถ้าคุณเป็นมือใหม่ แนะนำให้ศึกษาข้อมูลให้รอบคอบก่อน และมั่นใจว่าสำรองข้อมูลครบถ้วน ก่อนลงมืออ่านบทเรียนอย่างละเอียด ยืนยันรุ่นเครื่องและเวอร์ชันระบบของคุณ

\booksubsection{การรับรอง SafetyNet และ Play Integrity: ทำไมแอปธนาคารถึงใช้ไม่ได้หลังจาก Root}

ผู้ใช้ Root หลายคนเคยประสบกับช่วงเวลาสะดุดต่อไปนี้: หลังจาก Root แล้วตื่นเต้นที่จะติดตั้งโมดูลมากมาย แต่ผลคือเปิดแอปธนาคาร—“อุปกรณ์ของคุณไม่ปลอดภัย ไม่สามารถใช้งานได้” Google Pay ก็ใช้ไม่ได้ แม้แต่ Netflix ก็ปฏิเสธการเล่นวิดีโอความคมชัดสูง เบื้องหลังนี้คือกลไกการรับรองความสมบูรณ์ของอุปกรณ์ 2 ชุดของ Google: คือ SafetyNet รุ่นเก่าและ Play Integrity API รุ่นใหม่

\subsubsection*{SafetyNet / Play Integrity คืออะไร:}

\begin{itemize}
  \item \textbf{SafetyNet Attestation (รุ่นเก่า, เริ่มใช้งานปี 2017)}: API ตรวจสอบความสมบูรณ์ของอุปกรณ์ที่ Google เปิดตัวในช่วงแรก แอปจะใช้มันในการสอบถามไปยังเซิร์ฟเวอร์ของ Google ว่า "อุปกรณ์นี้ผ่านการทดสอบความเข้ากันได้แบบ CTS หรือไม่ และอยู่ในสถานะทางการที่ไม่ถูกดัดแปลงหรือไม่" ผลลัพธ์ที่ส่งกลับมาแบ่งเป็น Basic Attestation (การตรวจสอบพื้นฐาน) และ Strong Attestation (การตรวจสอบขั้นเข้มงวด ซึ่งต้องอาศัยคีย์ฮาร์ดแวร์) Google ได้ประกาศแล้วว่า SafetyNet จะค่อยๆ ยกเลิกการใช้งานหลังปี 2024
  \item \textbf{Play Integrity API (รุ่นใหม่, เริ่มใช้งานปี 2022)}: ทางเลือกที่ Google เปิดตัวขึ้นมาทดแทน ซึ่งเข้มงวดและครอบคลุมมากกว่า SafetyNet นอกเหนือจากความสมบูรณ์ของอุปกรณ์แล้ว ยังรวมถึง "ความสมบูรณ์ของแอป" (แอปถูกดัดแปลงหรือไม่) และ "ความสมบูรณ์ของบัญชี/การอนุญาต" (บัญชีผู้ใช้ได้รับความไว้วางใจจาก Google หรือไม่) แอปใหม่ควรย้ายไปใช้ Play Integrity
  \item \textbf{ใครใช้งานบ้าง}: แอปธนาคาร การเงิน การชำระเงิน (เช่น Google Pay, ธนาคารออนไลน์ต่างๆ, Alipay/WeChat Pay ในบางสถานการณ์), สตรีมมิ่งมีเดีย (การปกป้องลิขสิทธิ์ความคมชัดสูงของ Netflix), ระบบป้องกันการโกงในบางเกม (เช่นเกมบางส่วนของ Tencent/NetEase), แอปออฟฟิศขององค์กร (การตรวจสอบความสอดคล้องของอุปกรณ์ MDM)
\end{itemize}

\subsubsection*{ผลการรับรอง Play Integrity 3 ระดับ:}

\begin{itemize}
  \item \textbf{MEETS\_DEVICE\_INTEGRITY (ผ่านการตรวจสอบความสมบูรณ์ของอุปกรณ์)}: อุปกรณ์ทำงานบนระบบ Android แท้, Bootloader ไม่ได้ปลดล็อก, ไม่มีการ Root, ผ่านการทดสอบความเข้ากันได้ CTS สามารถใช้งานแอปธนาคารและแอปการชำระเงินได้ตามปกติ
  \item \textbf{MEETS\_BASIC\_INTEGRITY (ผ่านการตรวจสอบความสมบูรณ์พื้นฐาน)}: อุปกรณ์น่าเชื่อถือในระดับพื้นฐาน แต่อาจมีความผิดปกติเล็กน้อย แอปบางตัวอาจให้ใช้งานแบบลดทอนฟังก์ชัน (เช่น Netflix ยังเล่นได้แต่จำกัดความคมชัด) ในขณะที่แอปธนาคารที่เข้มงวดบางแอปจะปฏิเสธโดยตรง
  \item \textbf{MEETS\_STRONG\_INTEGRITY (ผ่านการตรวจสอบความสมบูรณ์ระดับเข้มงวด)}: ระดับสูงสุด ต้องมีการรับรองด้วยคีย์ฮาร์ดแวร์ (เช่น ชิป Titan M ในซีรีส์ Google Pixel) อุปกรณ์จากแบรนด์อื่นส่วนใหญ่แม้ไม่ได้ Root ก็ไม่สามารถทำระดับนี้ได้
  \item \textbf{ไม่มีป้ายกำกับใดๆ (NO\_INTEGRITY)}: อุปกรณ์ถูก Root, ปลดล็อก BL แล้ว, แฟลช ROM ของบุคคลที่สาม หรือถูกตรวจจับว่าเป็นโปรแกรมจำลอง แอปการชำระเงิน/ธนาคารส่วนใหญ่จะปฏิเสธการทำงานโดยตรง
\end{itemize}

\subsubsection*{ทำไมแอปเหล่านี้ถึงใช้งานไม่ได้หลังจาก Root:}

แอปจะส่งคำขอรับรองไปยังเซิร์ฟเวอร์ของ Google ผ่าน Play Integrity API จากนั้น Google จะตรวจสอบตัวชี้วัดหลายอย่างของอุปกรณ์—สถานะการล็อก Bootloader, พาร์ติชันระบบถูกแก้ไขหรือไม่, มีไฟล์ไบนารี su อยู่หรือไม่, SELinux อยู่ในโหมด permissive หรือไม่ เป็นต้น หากมีความผิดปกติแม้เพียงอย่างเดียว จะทำให้ระดับการรับรองที่ส่งกลับมาลดลงหรือว่างเปล่า เมื่อแอปธนาคารเห็นว่า “อุปกรณ์ไม่ปลอดภัย” ก็จะปฏิเสธการให้บริการตามธรรมชาติ

\subsubsection*{วิธีการซ่อนสมัยใหม่: Magisk Hide / Shamiko / Zygisk:}

เพื่อให้สามารถใช้งานแอปธนาคารไปพร้อมกับเพลิดเพลินกับความสะดวกสบายของ Root ชุมชนนักพัฒนาได้สร้างชุดวิธีการ “ซ่อนสถานะ Root” ขึ้นมาหลายวิธี แนวคิดหลักคือ: ทำให้โปรแกรมตรวจสอบเห็นว่าไม่ใช่ระบบที่ถูกดัดแปลง แต่เป็นอุปกรณ์ที่ “ดูเหมือนปกติ”

\begin{itemize}
  \item \textbf{รีเฟรชเครื่องมีความเสี่ยง, ดำเนินการต้องการระมัดระวัง อาจทำให้:}: วิธีการ Systemless Root หลักในปัจจุบัน—ไม่แก้ไขพาร์ติชันระบบ แต่ใช้วิธีการ mount เพื่อทำ Root ข้อดีของการออกแบบแบบนี้คือโอกาสที่ OTA จะอัปเดตล้มเหลวน้อย และการถอนการติดตั้งค่อนข้างสะอาด ถือเป็นพื้นฐานของวิธีการซ่อนทั้งหมดในภายหลัง แต่การใช้งาน Magisk ก็ยังคงต้องปลดล็อก BL อยู่ดี
  \item \textbf{Magisk Hide (หยุดพัฒนาแล้ว)}: โมดูลซ่อนที่มาพร้อมกับ Magisk ในช่วงแรก สามารถซ่อน Root จากแอปที่กำหนดได้ แต่ตั้งแต่ปี 2021 เป็นต้นมา ถูกลบออกจาก Magisk อย่างเป็นทางการแล้ว เนื่องจาก Google พัฒนาวิธีการตรวจจับอย่างต่อเนื่อง ทั้งสองฝ่ายจึงตกอยู่ในสถานการณ์ "เกมแมวจับหนู"
  \item \textbf{Zygisk}: กลไกใหม่ที่แนะนำใน Magisk 24.0+ คล้ายกับ Riru โดยทำการโหลดโมดูลผ่านการแทรกเข้าไปในโปรเซส Zygote ซึ่งซ่อนเร้นกว่าการแก้ไขระบบโดยตรง การเปิดใช้งาน Zygisk ถือเป็นข้อกำหนดเบื้องต้นสำหรับการใช้วิธีการซ่อนสมัยใหม่
  \item \textbf{Shamiko}: โมดูลซ่อนหลักในปัจจุบัน ทำงานร่วมกับ Zygisk สามารถซ่อน Root จากแอปที่กำหนด (โหมด "รายชื่อขาวแบบย้อนกลับ") ซ่อน Magisk ตัวมันเอง และปลอมแปลงสถานะการล็อก Bootloader โดย Shamiko จะถูกแจกจ่ายผ่านร้าน LSPosed หรือ GitHub
  \item \textbf{Play Integrity Fix(osm0sis)}: โมดูล Magisk ที่เจาะจงสำหรับการตรวจจับ Google Play Integrity โดยจะแก้ไขลายนิ้วมือของอุปกรณ์ (fingerprint) เพื่อให้ Google เข้าใจว่าอุปกรณ์ของคุณคือรุ่นทางการที่ไม่ได้ Root รุ่นใดรุ่นหนึ่ง
\end{itemize}

\begin{tcbwarn}{เกมแมวจับหนูจะไม่มีวันจบสิ้น}
Google จะอัปเดตกฎการตรวจจับเป็นระยะๆ และผู้ผลิตแอปก็จะตามไปอัปเกรดการตรวจสอบ Play Integrity เช่นกัน วิธีการซ่อนที่ใช้ได้ในวันนี้ อาจถูกตรวจจับในวันพรุ่งนี้ หากคุณพึ่งพาแอปธนาคาร/การชำระเงินอย่างเต็มที่ (เช่น โอนเงินประจำวัน, สแกนจ่ายเงิน) ทางเลือกที่ปลอดภัยที่สุดก็ยังคงเป็น: ไม่ต้อง Root, ไม่ต้องปลดล็อก BL วิธีการซ่อนเหมาะสำหรับผู้เล่นที่มีการตั้งค่า “เครื่องหลักไม่ Root + เครื่องสำรองเล่น Root” มากกว่า

\end{tcbwarn}

\booksubsection{คำแนะนำสำหรับผู้ใช้ทั่วไป}

\begin{itemize}
  \item \textbf{โหมดนักพัฒนา}: สามารถเปิดใช้งานเพื่อทดลองฟังก์ชันต่างๆ เช่น การเร่งความเร็วแอนิเมชัน, การวิเคราะห์การเรนเดอร์ GPU แต่ห้ามไปแก้ตัวเลือกอื่นๆ ที่ไม่เข้าใจโดยเด็ดขาด
  \item \textbf{การปลดล็อก OEM และ BL Lock}: หากคุณไม่ได้วางแผนจะแฟลชรอมหรือ Root ปล่อยให้สวิตช์ทั้งสองอยู่ในสถานะเริ่มต้น (ปิดอยู่) จะปลอดภัยที่สุด
  \item \textbf{ADB}: เครื่องมืออันทรงพลังที่ไม่ต้องใช้ Root ซึ่งคนทั่วไปก็ใช้ได้—เช่น ใช้\texttt{pm uninstall -\allowbreak -\allowbreak user 0}เพื่อถอนการติดตั้งแอปที่มาพร้อมเครื่อง แต่ก่อนลบต้องตรวจสอบให้แน่ใจก่อนว่าชื่อแพ็กเกจนั้นๆ สามารถลบได้หรือไม่
  \item \textbf{eSIM}: ไม่แนะนำให้ผู้ใช้ทั่วไปทำ Root เพราะความเสี่ยงมีมากกว่าผลตอบแทน หากต้องการใช้งานฟังก์ชันบางอย่างจริงๆ ควรตรวจสอบก่อนว่ามีทางเลือกที่ไม่ต้องใช้ Root หรือไม่ (เช่น ใช้ ADB ในการถอนการติดตั้งแอปที่มาพร้อมเครื่อง, หรือใช้ Shizuku แทนฟังก์ชัน Root บางส่วน)
\end{itemize}

\begin{tcbpractice}{ลงมือทำ: เปิดโหมดนักพัฒนาและปรับความเร็วแอนิเมชัน}
\begin{enumerate}
  \item : เข้าไปที่ การตั้งค่า > เกี่ยวกับโทรศัพท์
  \item : กด "หมายเลขรุ่น" ต่อเนื่อง 7 ครั้ง จนขึ้นข้อความว่าเปิดโหมดนักพัฒนาแล้ว
  \item : กลับไปที่การตั้งค่า หา "ตัวเลือกนักพัฒนา" (ปกติอยู่ท้ายการตั้งค่าหรือในหมวด "ระบบ")
  \item : หา "การปรับขนาดแอนิเมชันหน้าต่าง" "การปรับขนาดแอนิเมชันการเปลี่ยนผ่าน" "การปรับระยะเวลาแอนิเมชันโปรแกรม" แล้วตั้งเป็น 0.5x หรือปิดเสีย
  \item : ลองสัมผัสดูว่ามือถือเร็วขึ้นหรือไม่
\end{enumerate}

\end{tcbpractice}
